% Lösungen Einheit 4 von 5 – Vierecke nachweisen (zusammenbau v0.1)
\einheitenkopf{Einheit 4 von 5 · Vierecke nachweisen}

\erg{17}{a) Parallelogramm und ein rechter Winkel \quad b) $\overrightarrow{AB} = (3 | 1 | 2) = \overrightarrow{DC}$ (nicht mit $\overrightarrow{CD}$ vergleichen): $ABCD$ ist ein Parallelogramm. \quad c) $\overrightarrow{AB} = (2 | 2 | 1) = \overrightarrow{DC}$: Parallelogramm; $\overrightarrow{AB} \circ \overrightarrow{AD} = (2 | 2 | 1) \circ (-2 | 1 | 2) = 0$: rechter Winkel, also Rechteck. \quad d) $\overrightarrow{AB} = \overrightarrow{DC} = (2 | 2 | 1)$: Parallelogramm; $|\overrightarrow{AB}| = |\overrightarrow{AD}| = 3$: alle vier Seiten gleich lang (nicht nur zwei), also Raute. \quad e) $\overrightarrow{AB} = (3 | 1 | 0) = \overrightarrow{DC}$: Parallelogramm; $\overrightarrow{AB} \circ \overrightarrow{AD} = 5 \ne 0$: kein rechter Winkel bei $A$, also kein Rechteck. \quad f) $\overrightarrow{EF} = (-6 | 6 | 0) = 2 \cdot \overrightarrow{AB}$: $AB \parallel EF$, also Trapez (nicht aus gleichen Längen schließen); $|\overrightarrow{AE}| = |\overrightarrow{BF}| = \sqrt{18}$. \quad g) $\overrightarrow{AD} = (0 | 6 | 0) = 3 \cdot \overrightarrow{BC}$: Trapez. Höhe ist der Abstand der Parallelen in $x_1$-Richtung, $h = 6$ (nicht $|AB|$). Flächeninhalt $\frac12 \cdot (6 + 2) \cdot 6 = 24$ m². \quad h) $|\overrightarrow{KL}| = |\overrightarrow{KN}| = \sqrt{13}$ und $|\overrightarrow{ML}| = |\overrightarrow{MN}| = 5$: zwei Paare benachbarter gleich langer Seiten (die Gegenseiten sind verschieden lang). \quad i) $\overrightarrow{AE} = (0 | 6 | 1) = 2 \cdot \overrightarrow{CD}$ mit $\overrightarrow{CD} = (0 | 3 | 0{,}5)$: parallel; $\overrightarrow{CD} \circ \overrightarrow{DE} = (0 | 3 | 0{,}5) \circ (-6 | 0 | 0) = 0$: rechter Winkel bei $D$. \quad j) Die Seite an $P$ liegt in der Ebene und steht senkrecht auf $\overrightarrow{PQ} = (0 | 0 | 5)$: Richtung $(4 | -3 | 0)$, ihre Länge $5$ passt schon. $R = P + (4 | -3 | 0) = (6 | -2 | 0)$ (oder $(-2 | 4 | 0)$). \quad k) Diagonalenschnittpunkt ist der Spurpunkt der Geraden: $M(0 | 3 | 2)$. $\overrightarrow{MP} = (0 | 3 | 1)$, $\overrightarrow{MQ} = (0 | -1 | 3)$: gleich lang ($\sqrt{10}$), also ist $Q$ eine Ecke; $\overrightarrow{MP} \circ \overrightarrow{MQ} = 0$, also nicht die Gegenecke, sondern eine Nachbarecke.}
\erg{18}{a) Fehler: Gegenüber von $\overrightarrow{AB}$ liegt $\overrightarrow{DC}$, nicht $\overrightarrow{CD}$. Richtig: $\overrightarrow{DC} = (3 | 1 | 2) = \overrightarrow{AB}$, also Parallelogramm. \quad b) Ein Rechteck hat an jeder Ecke einen rechten Winkel; fehlt er an einer Ecke, ist die Bedingung schon verletzt. Zum Ausschließen genügt ein Gegenmerkmal, zum Nachweisen braucht man alle Merkmale.}
